\documentclass[main.tex]{subfiles}
\begin{document}
\section{Skia}
\label{sec:skia}

Skia is a widely used 2D graphics library.
Skia was originally developed by an independent company founded by
  engineers that worked on QuickDraw wich was the 2D graphics library and a core
  part of the classic MacOS.
%https://mergr.com/transaction/google-acquires-skia
The company behind Skia was quickly acquired by Google after its founding,
  with Skia becoming a core part of the Chromium browser engine's rendering
  stack from its inception.
Skia is actively maintained:
  as of March 11, 2026, its public main branch log shows 80 commits
  in the prior 7 days.
Its contributor base is also broad,
  with 109 entries in the AUTHORS file,
  and its release process remains regular,
  with a stable milestone cut every four weeks alongside Chromium.

Skia's basic APIs
  create drawable 2D ``surfaces'' and ``layers'',
  draw shapes to those surfaces and layers,
  blend them together,
  and apply effects.
These
  operations are issued through methods on a \skcanvas,
  which manages the drawing state.
At execution time, the behavior is backend-dependent:
  CPU backends apply draws
  directly to CPU-visible pixel memory,
  while GPU backends enqueue work that is flushed to backend API objects,
  and then submitted for GPU execution.
But in either case,
  this list of commands issued prior to rendering
  can be stored in an intermeddiate buffer of commands
  to be use for either serialization, replaying, or optimizations,
  which we can take to be a stand-alone programming language.
This section introduces the language's basic operations.

\subsection{Drawing and Painting}
\label{subsec:drawing-painting}

The \skdrawrect command is used to draw rectangles:%

\begin{figure}[H]
  \centering
  \begin{tabular}{c}
    \frame{
\includegraphics[scale=0.15]{drawing0.pdf}}
  \end{tabular}
  \quad
  \begin{tabular}{>{\smaller}T}
    SkPaint p;\\
    canvas->drawRect(rect1, p);
  \end{tabular}
  \quad
  \begin{tabular}{c}
    \frame{
\includegraphics[scale=0.15]{drawing1.pdf}}
  \end{tabular}
\end{figure}

% paint.setColor(SK_ColorRED);
% canvas->drawRect(SkRect::MakeXYWH(0, 0, 180, 180), paint);

\noindent
Here the layout is meant to evoke
  a Hoare triple:
the left square shows the (empty) state of the canvas
  before the \skdrawrect command,
  while the right square shows the state of the canvas afterwards.
The \skdrawrect command needs the bounds of the rectangle
  that needs to be drawn.
\skdrawrect also takes a \skpaint
  as an argument.
A \skpaint describes \emph{how} a geometry is drawn,
  for example how to \emph{fill} the geometry.
The default fill for a \skpaint is the color black,
  but we can also change it to another color.

\begin{figure}[H]
  \centering
  \begin{tabular}{c}
    \frame{
\includegraphics[scale=0.15]{drawing0.pdf}}
  \end{tabular}
  \quad
  \begin{tabular}{>{\smaller}T}
    SkPaint p;\\
    p.setColor(SK\_ColorRED);\\
    canvas->drawRect(rect1, p);
  \end{tabular}
  \quad
  \begin{tabular}{c}
    \frame{
\includegraphics[scale=0.15]{drawing2.pdf}}
  \end{tabular}
\end{figure}

\noindent
It's also possible to specify a non-constant fill,
  like a gradient:

\begin{figure}[H]
  \centering
  \begin{tabular}{c}
    \frame{
\includegraphics[scale=0.15]{drawing0.pdf}}
  \end{tabular}
  \quad
  \begin{tabular}{>{\smaller}T}
    SkPaint p;\\
    p.setShader(\\
    \quad SkShaders::LinearGradient(...)\\
    );\\
    canvas->drawRect(rect1, p);
  \end{tabular}
  \quad
  \begin{tabular}{c}
    \frame{
\includegraphics[scale=0.15]{drawing3.pdf}}
  \end{tabular}
\end{figure}

\noindent
Skia also supports complex fills like images and shaders.
  Formally, filling is a stage of Skia's drawing pipeline
  that assigns a color to every point in the shape.

\subsection{Blending}
\label{subsec:blending}

Things get a little more interesting
  once there is more than one shape being drawn.
Suppose we draw a second, overlapping, blue square:

% void draw(SkCanvas* canvas) {
% canvas->clear(SK_ColorWHITE);

% SkPaint paint;

% SkPoint pts[] = { {0, 0}, {180, 180} };

% SkColor4f colors[] = {
% SkColors::kBlue,
% SkColors::kRed
% };

%   SkScalar positions[] = { 0.0f, 1.0f };  // REQUIRED in your build

%   auto shader = SkShaders::LinearGradient(
%   pts,
%   { {colors, positions, SkTileMode::kClamp}, {} },
%   nullptr
%   );

%   paint.setShader(shader);

%   canvas->drawRect(SkRect::MakeXYWH(0, 0, 180, 180), paint);
% }

\begin{figure}[H]
  \centering
  \begin{tabular}{c}
    \frame{
\includegraphics[scale=0.15]{drawing2.pdf}}
  \end{tabular}
  \quad
  \begin{tabular}{>{\smaller}T}
    p.setColor(SK\_ColorBLUE);\\
    canvas->drawRect(rect2, p);
  \end{tabular}
  \quad
  \begin{tabular}{c}
    \frame{
\includegraphics[scale=0.15]{drawing4.pdf}}
  \end{tabular}
\end{figure}

% paint.setColor(SK_ColorBLUE);
% canvas->drawRect(SkRect::MakeXYWH(120, 120, 136, 136), paint);

In Skia, the process of drawing this second shape
  involves, conceptually, two stages.
First, as before, every pixel in \texttt{rect2}
  is assigned a fill color, here a solid blue.
Second, this assigned fill color is \emph{blended}
  with the existing color for that pixel.
The default \emph{blend mode}, called \sksrcover,
  just chooses the new color,%
\footnote{In reality, the blending is more complicated
  and depends on the opacity or \emph{alpha}
  of the new color,
  which can matter for anti-aliasing,
  but in our examples all the colors are opaque
  and we ignore anti-aliasing as a simplification.}
  but a program can select a different blend modes instead.
For example, the \skmultiply blend mode %
  combines both colors being blended:

\begin{figure}[H]
  \centering
  \begin{tabular}{c}
    \frame{
\includegraphics[scale=0.15]{drawing2.pdf}}
  \end{tabular}
  \quad
  \begin{tabular}{>{\smaller}T}
    p.setColor(SK\_ColorBLUE);\\
    p.setBlendMode(\\
    \quad SkBlendMode::kMultiply\\
    );\\
    canvas->drawRect(rect2, p);
  \end{tabular}
  \quad
  \begin{tabular}{c}
    \frame{\includegraphics[scale=0.15]{drawing5.pdf}}
  \end{tabular}
\end{figure}

Now, in the intersection of the two rectangles,
  the initial red color is blended with the newly-drawn blue fill,
  which in the \skmultiply blend mode results in black.
Formally, blending is another stage in the Skia drawing pipeline,
  which comes after filling and combines
  the fill color assigned to each pixel
  with the pre-existing color of that pixel.

\subsection{Clipping}
\label{subsec:clipping}

So far, \skdrawrect has been applying the fill
  to the entirety of the shape being drawn.
But, sometimes,
  we need more granual control over the drawing bounds.
Skia provides a dedicated \skcliprect command
  to restrict or \emph{clip} shapes to rectangles.%
\footnote{
  Skia allows clipping with arbitrary shapes,
  but the GPU back-end has fast paths for rectangles
  and rounded-rectangles, reducing them to scissor tests}
One uses \skcliprect like this:

\begin{figure}[H]
  \centering
  \begin{tabular}{c}
    \frame{
\includegraphics[scale=0.15]{drawing0.pdf}}
  \end{tabular}
  \quad
  \begin{tabular}{>{\smaller}T}
    canvas->clipRect(clip\_rect);\\
    canvas->drawRect(rect1, redPaint);
  \end{tabular}
  \quad
  \begin{tabular}{c}
    \frame{\includegraphics[scale=0.15]{drawing6.pdf}}
  \end{tabular}
\end{figure}

\noindent
Note the unusual order of operations here.
The \skcliprect command is run first,
  and it does not modify the image at all.
Instead, it modifies the shape drawn
  by \emph{later} draw operations.
Formally, clipping is yet another stage of the Skia drawing pipeline;
  it comes before filling and modifies the shape being drawn.

Semantically, \skcliprect modifies the \emph{clip state},
  which stored in the canvas object and affects drawing commands.
This means that a single \skcliprect operation
  can affect multiple \skdrawrect operations:

\begin{figure}[H]
  \centering
  \begin{tabular}{c}
    \frame{
\includegraphics[scale=0.15]{drawing0.pdf}}
  \end{tabular}
  \quad
  \begin{tabular}{>{\smaller}T}
    canvas->clipRect(c\_rect);\\
    canvas->drawRect(rect1, redPaint);\\
    canvas->drawRect(rect2, bluePaint);
  \end{tabular}
  \quad
  \begin{tabular}{c}
    \frame{
\includegraphics[scale=0.15]{drawing7.pdf}}
  \end{tabular}
\end{figure}

Naturally, one might want to clip
  only a subset of drawing commands.
Luckily, Skia maintains a \emph{stack} of clip states;
  \skcliprect affects the top of the stack,
  which is used by drawing commands,
  while \sksave and \skrestore commands
  push and pop from the stack.
  \sksave and \skrestore act like brackets,
  restricting the effect of \skcliprect to the enclosed commands:

\begin{figure}[H]
  \centering
  \begin{tabular}{c}
    \frame{
\includegraphics[scale=0.15]{drawing0.pdf}}
  \end{tabular}
  \quad
  \begin{tabular}{>{\smaller}T}
    canvas->save();\\
    canvas->clipRect(c\_rect);\\
    canvas->drawRect(rect1, redPaint);\\
    canvas->restore();\\
    canvas->drawRect(rect2, bluePaint);
  \end{tabular}
  \quad
  \begin{tabular}{c}
    \frame{\includegraphics[scale=0.15]{drawing8.pdf}}
  \end{tabular}
\end{figure}

\subsection{Layers}
\label{subsec:layers}

Much like \sksave and \skrestore can save clip state,
  there is also a \sksavelayer and \skrestore command
  that saves and restores the drawing surface itself.
Conceptually, \sksavelayer allocates
  a new, empty drawing surface.

\begin{figure}[H]
  \centering
  \begin{tabular}{c}
    \frame{
\includegraphics[scale=0.15]{drawing2.pdf}}
  \end{tabular}
  \quad
  \begin{tabular}{>{\smaller}T}
    canvas->saveLayer();
  \end{tabular}
  \quad
  \begin{tabular}{c}
    \frame{
\includegraphics[scale=0.15]{drawing0.pdf}}\\[0.6em]
    \frame{
\includegraphics[scale=0.15]{drawing2.pdf}}
  \end{tabular}
\end{figure}

All drawing commands until the corresponding \skrestore
  draw to this new surface.

\begin{figure}[H]
  \centering
  \begin{tabular}{c}
    \frame{
\includegraphics[scale=0.15]{drawing0.pdf}}\\[0.6em]
    \frame{
\includegraphics[scale=0.15]{drawing2.pdf}}
  \end{tabular}
  \quad
  \begin{tabular}{>{\smaller}T}
    canvas->drawRect(rect2, bluePaint);
  \end{tabular}
  \quad
  \begin{tabular}{c}
    \frame{
\includegraphics[scale=0.15]{drawing9.pdf}}\\[0.6em]
    \frame{
\includegraphics[scale=0.15]{drawing2.pdf}}
  \end{tabular}
\end{figure}

Then, when the \skrestore command it reached,
  the whole resulting drawing surface is blended back into
  the original drawing surface.

\begin{figure}[H]
  \centering
  \begin{tabular}{c}
    \frame{
\includegraphics[scale=0.15]{drawing9.pdf}}\\[0.6em]
    \frame{
\includegraphics[scale=0.15]{drawing2.pdf}}
  \end{tabular}
  \quad
  \begin{tabular}{>{\smaller}T}
    restore();
  \end{tabular}
  \quad
  \begin{tabular}{c}
    \frame{
\includegraphics[scale=0.15]{drawing4.pdf}}
  \end{tabular}
\end{figure}

Semantically, you can think of \sksavelayer commands
  as re-associating the order of blend operations.
In the \sksave example,
  the order of operations is

\begin{equation*}
  \textsf{blend}(\textsf{blend}(\textsf{empty}, \textsf{red}), \textsf{blue}),
\end{equation*}

While in the \sksavelayer example
  the order of operations is

\begin{equation*}
  \textsf{blend}(\textsf{blend}(\textsf{empty}, \textsf{red}), \textsf{blend}(\textsf{empty}, \textsf{blue})).
\end{equation*}

With \sksrcover blending, which is associative,
  this order of operations is not important,
  but with other blend modes it becomes significant.
For example,
  we reepat the previous snippet,
  but pass a \skmultiply blend mode to the \sksavelayer command.

\begin{figure}[H]
  \centering
  \begin{tabular}{c}
    \frame{
\includegraphics[scale=0.15]{drawing9.pdf}}\\[0.6em]
    \frame{
\includegraphics[scale=0.15]{drawing2.pdf}}
  \end{tabular}
  \quad
  \begin{tabular}{>{\smaller}T}
    restore(); // multiply blending
  \end{tabular}
  \quad
  \begin{tabular}{c}
    \frame{\includegraphics[scale=0.15]{drawing5.pdf}}
  \end{tabular}
\end{figure}

\subsubsection{Masking}
\label{subsubsec:masking}

To further motivate the need for \sksavelayer,
  consider masking again.
In \Cref{subsec:clipping} we saw how we can use
  \skcliprect to mask images.
This works when we only need to mask by a rectangle.
But masking with complicated shapes is hard with just clipping.
Instead of trying to clip directly with a complicated shape,
  we can instead draw the clip shape in stages on a seperate layer
  and mask using that shape.

\begin{figure}[H]
  \centering
  \begin{tabular}{c}
    \frame{\includegraphics[scale=0.15]{mandrill.pdf}}
  \end{tabular}
  \quad
  \begin{tabular}{>{\smaller}T}
    SkPaint blend\_p;\\
    blend\_p.setBlendMode(\\
    \quad SkBlendMode::kDstIn\\
    );\\
    canvas->saveLayer(blendPaint);
  \end{tabular}
  \quad
  \begin{tabular}{c}
    \frame{
\includegraphics[scale=0.15]{drawing0.pdf}}\\[0.6em]
    \frame{\includegraphics[scale=0.15]{mandrill.pdf}}
  \end{tabular}
\end{figure}

We now draw the mask shape on the new layer.

\begin{figure}[H]
  \centering
  \begin{tabular}{c}
    \frame{
\includegraphics[scale=0.15]{drawing0.pdf}}\\[0.6em]
    \frame{\includegraphics[scale=0.15]{mandrill.pdf}}
  \end{tabular}
  \quad
  \begin{tabular}{>{\smaller}T}
    SkPaint p;\\
    canvas->drawRect(rect1, p);\\
    canvas->drawRect(rect2, p);
  \end{tabular}
  \quad
  \begin{tabular}{c}
    \frame{
\includegraphics[scale=0.15]{mask.pdf}}\\[0.6em]
    \frame{\includegraphics[scale=0.15]{mandrill.pdf}}
  \end{tabular}
\end{figure}

\skdstin emulates masking through blending;
When we \skdstin one pixel over another,
  an opaque top pixel keeps the bottom pixel,
  and a transparent top pixel removes it.

\begin{figure}[H]
  \centering
  \begin{tabular}{c}
    \frame{
\includegraphics[scale=0.15]{mask.pdf}}\\[0.6em]
    \frame{\includegraphics[scale=0.15]{mandrill.pdf}}
  \end{tabular}
  \quad
  \begin{tabular}{>{\smaller}T}
    restore();
  \end{tabular}
  \quad
  \begin{tabular}{c}
    \frame{\includegraphics[scale=0.15]{masked_mandrill.png}}
  \end{tabular}
\end{figure}

\skdstin blending is also far more powerful than just clipping,
  because it also allows masking with gradients,
  allowing us to create more complicated images
  that just clipping would not allow us to.
For example,
  here is the previous image,
  now masked with a radial gradient.

\begin{figure}[H]
  \centering
  \begin{tabular}{c}
    \frame{
\includegraphics[scale=0.15]{gradient.png}}\\[0.6em]
    \frame{\includegraphics[scale=0.15]{mandrill.pdf}}
  \end{tabular}
  \quad
  \begin{tabular}{>{\smaller}T}
    restore();
  \end{tabular}
  \quad
  \begin{tabular}{c}
    \frame{\includegraphics[scale=0.15]{gradient_mandrill.png}}
  \end{tabular}
\end{figure}

\subsection{Equivalences}
\label{subsec:equivalences}

In \Cref{subsubsec:masking} we saw how complex masks
  can be done with ease using a \sksavelayer
  with \skdstin blending.
But what if just a rectangle mask is needed?
Our \sksavelayer and \skdstin program would still work.
But,
  it is clearly less performant than just doing a clip;
We have to allocate a whole new layer,
  just to do something Skia already has a fast path for!
Here,
  we have two different Skia programs that produce the same image,
  but one is clearly a better choice than the other.

\begin{figure}[H]
  \centering
  \begin{tabular}{c}
    \frame{
\includegraphics[scale=0.15]{drawing0.pdf}}
  \end{tabular}
  \quad
  \begin{tabular}{l}
    \begin{tabular}{>{\smaller}T}
      canvas->clipRect(rect1);\\
      canvas->drawImage(...);
    \end{tabular}\\
    \noalign{\vskip 1.4em}
    \cline{1-1}
    \noalign{\vskip 1.4em}
    \begin{tabular}{>{\smaller}T}
      canvas->drawImage(...);\\
      SkPaint b;\\
      b.setBlendMode(SkBlendMode::kDstIn);\\
      canvas->saveLayer(b);\\
      SkPaint p;\\
      canvas->drawRect(rect1, p);\\
      canvas->restore();
    \end{tabular}
  \end{tabular}
  \quad
  \begin{tabular}{c}
    \frame{\includegraphics[scale=0.15]{rect_mandrill.png}}
  \end{tabular}
\end{figure}

Therefore it should be possible to replace a Skia program
  with a faster one,
  given they have identical results.
In \Cref{sec:semantics}
  we develop a semantics,
  which allows proving that two Skia programs have identical results,
  including verifying general-purpose rewrite rules.
\Cref{sec:compiler} then uses these rewrites
  to build a general-purpose Skia optimizer
  that rewrites programs to reduce the use
  of multiple layers.
\Cref{sec:eval} then shows that this compiler
  results in dramatically faster rendering
  on a wide range of real-world Skia programs.

\end{document}
